\documentclass{article}
\usepackage[T1]{fontenc}
\usepackage[header=false,clearsection=false]{../../../adenc}

\newcommand{\answer}[1]{\par\medskip #1\par\medskip}
\setlist[enumerate,1]{label=(\arabic*),leftmargin=*}

\begin{document}
Consider the SCM with \vocab{additive unobservable confounder}
\begin{equation*}
    Y = f(A, X) + f_U(U),
\end{equation*}
where there is no interaction between \((A, X)\) and \(U\), the unobserved confounder. Then,
\begin{equation*}
    \E(Y\mid X, \doo(A = a)) - \E(Y\mid  X, \doo(A = a')) = f(a, X)- f(a', X).
\end{equation*}
Our goal is to learn the contrast function \(f\).

\begin{theorem}
    \label{thm:iv}
    Under IV assumptions and additive unobservable confounders, there exists \(\tilde{f}\) such that
    \begin{enumerate}
        \item \(\E(Y\mid X, \doo(A = a)) -\E(Y\mid X, \doo(A = a')) = \tilde{f}(a, x) - \tilde{f}(a', x)\).
        \item \(\E(Y \mid Z, X) = \int\tilde{f}(a, x)p(a\mid z, x)\dd{a}\), where \(p(a\mid z, x)\) is the probability density of treatment.
    \end{enumerate}
    In particular, if there exists a unique function \(g\) satisfying condition 2, then \(\tilde{f} = g\).
\end{theorem}

\begin{proof}
    By unconfoundedness, \(U\perp Z\mid X\). Then,
    \begin{align*}
        \E(Y\mid Z, X) & = \E(f(A, X)\mid Z, X) + \E(f_U(U)\mid Z, X)\\
        & = \E(f(A, X)\mid Z, X) + \underbrace{\E(f_U(U)\mid X)}_{h(X)}\\
        & = \E(\tilde{f}(A, X)\mid Z, X), \text{ where }\tilde{f}(A, X)\coloneqq f(A, X) + h(X)\\
        & = \int\tilde{f}(a, X)p(a\mid Z, X)\dd{a}.
    \end{align*}
\end{proof}
Note that we never used relevance in the proof. However, we need it for the theorem to be useful. The idea is that the theorem \Cref{thm:iv} gives a system of equations of the form
\begin{equation*}
    \E(Y\mid \cdot) = \langle\tilde{f}(\cdot,X ), p(\cdot, z)\rangle
\end{equation*}
If the instrument is strong, we have multiply conditional distribution of the form \(p(\cdot, z)\) by varying \(z\), which narrows the choice for function \(\tilde{f}\). This means that weaker instruments yields higher fundamental uncertainty even with infinite data.

\subsubsection{Model fitting}
\Cref{thm:iv} states that we can fit the model in two steps:
\begin{enumerate}
    \item Fit a conditional probability \(\hat{p}(a\mid x, z)\);
    \item Solve the minimization problem
    \begin{equation*}
        \tilde{f} = \argmin_f \frac{1}{n}\sum_{i=1}^n\left(Y - \int \hat{p}(a\mid x, z)f(a, x)\dd{a}\right)^2
    \end{equation*}
    \item Compute the average local treatment effect estimate
    \begin{equation*}
        \frac{1}{n}\sum_{i = 1}^n\frac{\partial}{\partial a}\tilde f(A_i, X_i)
    \end{equation*}
\end{enumerate}
